% =====================================================================
% PROPOSED MANUSCRIPT ADDITIONS — exploratory moderator analyses
% Drafted to accompany the online-supplement chapter "moderator_analysis".
% These are SNIPPETS to paste into main.tex; they are not compiled here.
% All numbers are taken from the rendered supplement chapter.
% =====================================================================


% ---------------------------------------------------------------------
% (1) RESULTS — new subsection.
% Place immediately AFTER the existing subsection
% "Uni-moderator analyses of taxonomic group".
% ---------------------------------------------------------------------
\subsection{Exploratory uni-moderator analyses of extracted moderators}

Beyond taxonomic group, we extracted four candidate moderators during data
collection: the copying mechanism (generalized \textit{vs.} individual), the
modality of the social cue (visual, chemical, or acoustic), the virginity of the
observer, and the relative age of the demonstrator. These were not part of our
pre-registered analyses and several of their levels are sparsely sampled, so we
treat them as exploratory and report them in full in the online supplement.
The relative age of the demonstrator showed no variation in our extraction (all
``same-age'') and could not be analyzed. The remaining three moderators were fitted
as single-moderator multilevel models on the new extraction (both scales) and,
where the coding could be harmonized with an original dataset, on the combined
data. No moderator explained a detectable share of the heterogeneity on either
scale (all Wald $Q_M$, $p > 0.2$; total $I^2$ remained above 80\%). One contrast
is nonetheless worth noting: whereas Jones and DuVal~\cite{jonesMechanismsSocialInfluence2019}
reported stronger copying among virgin observers, our new extraction showed the
opposite ordering (Hedges' \textit{g} = 0.09 for virgin \textit{vs.} 0.36 for
non-virgin observers; Wald $Q_M$ = 1.61, $p$ = 0.22), and the difference all but
disappeared when the new data were pooled with the original binary data
(OR = 1.75 \textit{vs.} 1.91; $p$ = 0.72). The contrast was not statistically
distinct on either scale. Full diagnostics, model summaries, between-level
contrasts, and orchard plots are provided in the online supplement.


% ---------------------------------------------------------------------
% (2) DISCUSSION — replacement for the candidate-moderators sentence in the
% "Knowledge gaps and future opportunities" subsection.
%
% REPLACE the existing sentence:
%   "We collected several candidate moderators during extraction---the copying
%    mechanism (generalized versus individual), the modality of the social cue
%    (visual, chemical, or acoustic), the virginity of the observer, and the
%    relative age of the demonstrator---and we encourage future syntheses,
%    ideally pre-registered, to test these explicitly as the dataset continues
%    to grow."
% WITH the following:
% ---------------------------------------------------------------------
We collected several candidate moderators during extraction---the copying
mechanism (generalized versus individual), the modality of the social cue
(visual, chemical, or acoustic), the virginity of the observer, and the relative
age of the demonstrator---and explored them in post-hoc, non-pre-registered
uni-moderator analyses (online supplement). None explained a detectable share of
the heterogeneity, which is consistent with the variation residing at
finer-grained or interactive scales than our coding captured, and with several
levels (the chemical and acoustic cues in particular) being too sparsely sampled
to support inference. One result nonetheless deserves emphasis: the updated
evidence did not reproduce---and, in the new studies, tended to reverse---the
stronger copying among virgin observers reported by Jones and
DuVal~\cite{jonesMechanismsSocialInfluence2019}. We therefore encourage future
syntheses, ideally pre-registered, to test these moderators explicitly as the
dataset continues to grow, treating even the conditionality signals inherited
from the original meta-analyses as hypotheses to be re-examined rather than
settled results.


% ---------------------------------------------------------------------
% (3) DISCUSSION — OPTIONAL caveat in the social-learning subsection.
% In the paragraph beginning "Theory does not predict that animals should copy
% indiscriminately...", the sentence currently reads:
%   "Jones and DuVal found exactly such conditionality, with stronger copying
%    among inexperienced (virgin) females, in free-living settings, and when
%    social information ran counter to an innate preference;"
% Consider appending the parenthetical caveat shown here so the cited virginity
% effect is not presented as uncontested by our own update:
% ---------------------------------------------------------------------
Jones and DuVal~\cite{jonesMechanismsSocialInfluence2019} found exactly such
conditionality, with stronger copying among inexperienced (virgin) females
(although, as noted below and in the online supplement, our updated data did not
reproduce this particular virginity effect), in free-living settings, and when
social information ran counter to an innate preference;
